% GENERATED FILE. Edit canonical lessons, then run npm run generate:books.
% canonical-source-hash: fnv1a64:eeaffbc0d021f0cd
% canonical-lessons: AR-C472-sakhra, AR-C472-hadba, AR-C472-qimma, AR-C472-munhadar, AR-C472-haffa

\chapter{Rocks and peaks}
\label{ch:ar-rocks-and-peaks}
\hlchaptermodality{472}

\begin{chapteropening}
\textbf{By the end of this chapter:} \emph{I can say a rock, a plateau, a summit, a slope and an edge.}

Complete the last lesson of chapter 472: I can say a rock, a plateau, a summit, a slope and an edge.
\end{chapteropening}

\section[\texorpdfstring{\ar{صخرة} (ṣakhra)}{صخرة (ṣakhra)}]{\ar{صخرة} (ṣakhra) — a rock}

\label{lesson:AR-C472-sakhra}



\begin{quote}
Before the new one: say the Arabic for a current, then the Arabic for depth.
\end{quote}

\subsection*{You'll want to know: \ar{صخرة}}
\textbf{\ar{صخرة}} — \emph{ṣakhra} — \textquotedblleft{}a rock\textquotedblright{}.

\begin{quote}
\ar{جلسنا} \ar{على} \ar{صخرة} \ar{كبيرة}.
\end{quote}

\begin{quote}
\emph{jalasnā ʿalā ṣakhra kabīra.}
\end{quote}

\begin{quote}
We sat on a big rock.
\end{quote}

Plural: \ar{صخور}, \emph{ṣukhūr}.

\begin{cousinweb}
\textbf{\ar{صخرة}} (\emph{ṣakhra}) is built on the native Arabic root \emph{ṣ-kh-r}, \textquotedblleft{}great hard rocks\textquotedblright{} (Lane).
\end{cousinweb}

\begin{grammarlens}[title={What you've built}]
One more word for getting around.
\end{grammarlens}

\subsection*{Your turn}
\begin{itemize}
\raggedright
  \item \emph{Say it:} \emph{ṣakhra}
  \item \emph{Say it:} \emph{ṣakhra}, once more
  \item \emph{Say it:} \emph{jalasnā ʿalā ṣakhra kabīra.} — We sat on a big rock.
  \item \emph{Recall:} say the Arabic for a current, then the Arabic for depth, then say \emph{ṣakhra} again
\end{itemize}

\begin{tcolorbox}[breakable,colback=teal!4,colframe=teal!35!black,title={Before you move on}]
What does \ar{صخرة} mean? (\textquotedblleft{}A rock\textquotedblright{}.) Say it once more.
\end{tcolorbox}
\section[\texorpdfstring{\ar{هضبة} (haḍba)}{هضبة (haḍba)}]{\ar{هضبة} (haḍba) — a plateau; a hill}

\label{lesson:AR-C472-hadba}



\begin{quote}
Before the new one: say the Arabic for depth, then the Arabic for a rock.
\end{quote}

\subsection*{You'll want to know: \ar{هضبة}}
\textbf{\ar{هضبة}} — \emph{haḍba} — \textquotedblleft{}a plateau; a hill\textquotedblright{}.

\begin{quote}
\ar{تقع} \ar{القرية} \ar{على} \ar{هضبة} \ar{واسعة}.
\end{quote}

\begin{quote}
\emph{taqaʿu al-qarya ʿalā haḍba wāsiʿa.}
\end{quote}

\begin{quote}
The village lies on a wide plateau.
\end{quote}

\begin{cousinweb}
\textbf{\ar{هضبة}} (\emph{haḍba}) is built on the native Arabic root \emph{h-ḍ-b}: Lane gives \emph{haḍba} as a hill, or a mountain spreading over the surface of the ground, beside the root's sense of steady rain (Lane). Modern Arabic uses it for a plateau.
\end{cousinweb}

\begin{grammarlens}[title={What you've built}]
One more word for getting around.
\end{grammarlens}

\subsection*{Your turn}
\begin{itemize}
\raggedright
  \item \emph{Say it:} \emph{haḍba}
  \item \emph{Say it:} \emph{haḍba}, once more
  \item \emph{Say it:} \emph{taqaʿu al-qarya ʿalā haḍba wāsiʿa.} — The village lies on a wide plateau.
  \item \emph{Recall:} say the Arabic for depth, then the Arabic for a rock, then say \emph{haḍba} again
\end{itemize}

\begin{tcolorbox}[breakable,colback=teal!4,colframe=teal!35!black,title={Before you move on}]
What does \ar{هضبة} mean? (\textquotedblleft{}A plateau; a hill\textquotedblright{}.) Say it once more.
\end{tcolorbox}
\section[\texorpdfstring{\ar{قمة} (qimma)}{قمة (qimma)}]{\ar{قمة} (qimma) — a summit; a peak}

\label{lesson:AR-C472-qimma}



\begin{quote}
Before the new one: say the Arabic for a rock, then the Arabic for a plateau.
\end{quote}

\subsection*{You'll want to know: \ar{قمة}}
\textbf{\ar{قمة}} — \emph{qimma} — \textquotedblleft{}a summit; a peak\textquotedblright{}.

\begin{quote}
\ar{وصلنا} \ar{إلى} \ar{قمة} \ar{الجبل} \ar{بعد} \ar{الظهر}.
\end{quote}

\begin{quote}
\emph{waṣalnā ilā qimmat al-jabal baʿda aẓ-ẓuhr.}
\end{quote}

\begin{quote}
We reached the summit of the mountain in the afternoon.
\end{quote}

Plural: \ar{قمم}, \emph{qimam}.

A meeting of heads of state is also a \emph{qimma}.

\begin{cousinweb}
\textbf{\ar{قمة}} (\emph{qimma}) is built on the native Arabic root \emph{q-m-m}, \textquotedblleft{}the top of a thing\textquotedblright{} (Lane).
\end{cousinweb}

\begin{grammarlens}[title={What you've built}]
One more word for getting around.
\end{grammarlens}

\subsection*{Your turn}
\begin{itemize}
\raggedright
  \item \emph{Say it:} \emph{qimma}
  \item \emph{Say it:} \emph{qimma}, once more
  \item \emph{Say it:} \emph{waṣalnā ilā qimmat al-jabal baʿda aẓ-ẓuhr.} — We reached the summit of the mountain in the afternoon.
  \item \emph{Recall:} say the Arabic for a rock, then the Arabic for a plateau, then say \emph{qimma} again
\end{itemize}

\begin{tcolorbox}[breakable,colback=teal!4,colframe=teal!35!black,title={Before you move on}]
What does \ar{قمة} mean? (\textquotedblleft{}A summit; a peak\textquotedblright{}.) Say it once more.
\end{tcolorbox}
\section[\texorpdfstring{\ar{منحدر} (munḥadar)}{منحدر (munḥadar)}]{\ar{منحدر} (munḥadar) — a slope}

\label{lesson:AR-C472-munhadar}



\begin{quote}
Before the new one: say the Arabic for a plateau, then the Arabic for a summit.
\end{quote}

\subsection*{You'll want to know: \ar{منحدر}}
\textbf{\ar{منحدر}} — \emph{munḥadar} — \textquotedblleft{}a slope\textquotedblright{}.

\begin{quote}
\ar{نزلنا} \ar{المنحدر} \ar{ببطء}.
\end{quote}

\begin{quote}
\emph{nazalnā al-munḥadar bi-buṭʾ.}
\end{quote}

\begin{quote}
We went down the slope slowly.
\end{quote}

\begin{cousinweb}
\textbf{\ar{منحدر}} (\emph{munḥadar}) is built on the native Arabic root \emph{ḥ-d-r}, \textquotedblleft{}to descend\textquotedblright{} (Lane). It is the place-noun of Form VII, \textquotedblleft{}the way down\textquotedblright{}.
\end{cousinweb}

\begin{grammarlens}[title={What you've built}]
One more word for getting around.
\end{grammarlens}

\subsection*{Your turn}
\begin{itemize}
\raggedright
  \item \emph{Say it:} \emph{munḥadar}
  \item \emph{Say it:} \emph{munḥadar}, once more
  \item \emph{Say it:} \emph{nazalnā al-munḥadar bi-buṭʾ.} — We went down the slope slowly.
  \item \emph{Recall:} say the Arabic for a plateau, then the Arabic for a summit, then say \emph{munḥadar} again
\end{itemize}

\begin{tcolorbox}[breakable,colback=teal!4,colframe=teal!35!black,title={Before you move on}]
What does \ar{منحدر} mean? (\textquotedblleft{}A slope\textquotedblright{}.) Say it once more.
\end{tcolorbox}
\section[\texorpdfstring{\ar{حافة} (ḥāffa)}{حافة (ḥāffa)}]{\ar{حافة} (ḥāffa) — an edge; a rim}

\label{lesson:AR-C472-haffa}



\begin{quote}
Before the new one: say the Arabic for a summit, then the Arabic for a slope.
\end{quote}

\subsection*{You'll want to know: \ar{حافة}}
\textbf{\ar{حافة}} — \emph{ḥāffa} — \textquotedblleft{}an edge; a rim\textquotedblright{}.

\begin{quote}
\ar{وقفنا} \ar{على} \ar{حافة} \ar{الوادي}.
\end{quote}

\begin{quote}
\emph{waqafnā ʿalā ḥāffat al-wādī.}
\end{quote}

\begin{quote}
We stood on the edge of the valley.
\end{quote}

\begin{cousinweb}
\textbf{\ar{حافة}} (\emph{ḥāffa}) is usually traced to the native Arabic root \emph{ḥ-f-f}, \textquotedblleft{}to surround, to border\textquotedblright{}.
\end{cousinweb}

\begin{grammarlens}[title={What you've built}]
One more word for getting around.
\end{grammarlens}

\subsection*{Your turn}
\begin{itemize}
\raggedright
  \item \emph{Say it:} \emph{ḥāffa}
  \item \emph{Say it:} \emph{ḥāffa}, once more
  \item \emph{Say it:} \emph{waqafnā ʿalā ḥāffat al-wādī.} — We stood on the edge of the valley.
  \item \emph{Recall:} say the Arabic for a summit, then the Arabic for a slope, then say \emph{ḥāffa} again
\end{itemize}

\begin{tcolorbox}[breakable,colback=teal!4,colframe=teal!35!black,title={Before you move on}]
What does \ar{حافة} mean? (\textquotedblleft{}An edge; a rim\textquotedblright{}.) Say it once more.
\end{tcolorbox}
